\chapter{Mineral reaction fronts}

This exercise introduces the simulation of a calcite dissolution front, illustrating one of the most
fundamental coupled hydrogeochemical processes: the displacement of a
water in chemical equilibrium with a mineral by a solution that is undersaturated with respect to this mineral.

The conceptual model consists of a one-dimensional column (100\,m long,
$1\,\mathrm{m} \times 1\,\mathrm{m}$ cross-section) initially saturated
with a water in equilibrium with calcite.  A sodium chloride solution
(\ce{NaCl}, $0.001\,\mathrm{mol/L}$, $\mathrm{pH} = 7$) is continuously
injected at the upstream boundary at a constant flow rate.  This inflow
water contains no calcium and no inorganic carbon, and is therefore
strongly undersaturated with respect to calcite.  As the \ce{NaCl} front
advances through the column, calcite dissolves wherever the injected
water contacts the mineral, releasing \ce{Ca^{2+}} and carbonate species
into solution.  A sharp dissolution front develops and propagates
downstream.

This exercise is relevant to a wide range of natural and engineered
systems, including the flushing of carbonate aquifers, \ce{CO2} storage,
mine site drainage, and \textit{in-situ} leaching of mineral ore deposits.

% ============================================================
\section{Spatial Discretisation and Flow Model}
% ============================================================

To implement this model setup in ORTI3D proceed as follows.  A summary
of all relevant parameters is provided in Table~\ref{tab:params}.

\begin{table}[ht]
\centering
\caption{Flow and solute transport parameters.}
\label{tab:params}
\begin{tabular}{@{}L{8.5cm}l@{}}
\toprule
\textbf{Parameter} & \textbf{Value} \\
\midrule
Flow simulation type              & Steady state \\
Total simulation time (days)      & 100          \\
Time step size (days)             & 1            \\
Number of columns (grid cells)    & 20           \\
Grid spacing $\Delta x$ (m)       & 5            \\
Model length (m)                  & 100          \\
Model width $\times$ height (m$\times$m) & $1 \times 1$ \\
Hydraulic conductivity (m/day)    & 10           \\
Effective porosity ($-$)          & 0.25         \\
Total porosity ($-$)              & 0.25         \\
Pore velocity (m/day)             & 2            \\
Injection flow rate (m$^3$/day)   & 0.5          \\
Longitudinal dispersivity $\alpha_L$ (m) & 0.01  \\
\bottomrule
\end{tabular}
\end{table}

\begin{itemize}

\item In ORTI3D, create a new model and save it in a new and separate
  directory to avoid any confusion with other models.

\item Under \opath{Parameters $\rightarrow$ Model} select \gui{Grid} and
  specify the model domain.  Define a one-dimensional column: set the
  number of columns to \gui{20}, the column width to \gui{5\,m}, 1 row
  of 1\,m width, and 1 layer of 1\,m height.  Under
  \opath{Parameters $\rightarrow$ Model} select \gui{Model} and set the
  Dimension to \gui{1D/Column} and the Type to \gui{Confined}.

\item Under \opath{Parameters $\rightarrow$ Time} set the \gui{Total
  Simulation Time} to \gui{100 days} and the \gui{Step Size} to
  \gui{1 day}.  Leave the Step Mode as Linear.

\item Go to \opath{Spatial Attributes $\rightarrow$ MODFLOW $\rightarrow$
  DIS $\rightarrow$ dis.6 Top of Layers (TOP)} and set the top of the
  model to \gui{1\,m}.  Confirm with OK\@.  The bottom defaults to 0\,m
  and does not need to be changed.

\item Set boundary conditions for hydraulic heads by going to
  \opath{Spatial Attributes $\rightarrow$ MODFLOW $\rightarrow$ BAS6
  $\rightarrow$ bas.3 Boundary Conditions}.  Set the background value to
  \gui{1} (active cell) for all cells.  Then select the Zone Tool and
  draw a point at the last (downstream) cell.  Set the value of this
  zone to \gui{$-1$} to define it as a fixed-head boundary condition.

\item Set initial hydraulic heads by going to \opath{Spatial Attributes
  $\rightarrow$ MODFLOW $\rightarrow$ BAS6 $\rightarrow$ bas.5 Initial
  Heads} and enter a background value of \gui{1\,m}.  The head at the
  fixed-head downstream boundary will remain at this value throughout
  the simulation.

\item Specify the hydraulic conductivity by going to \opath{Spatial
  Attributes $\rightarrow$ MODFLOW $\rightarrow$ LPF $\rightarrow$ lpf.8
  Horizontal Hydraulic Conductivity} and set the value to
  \gui{10\,m/day}.

\item Add an injection well at the upstream (first) cell by going to
  \opath{Spatial Attributes $\rightarrow$ MODFLOW $\rightarrow$ WEL} and
  creating a point Zone at the first cell.  Set the injection flow rate
  to \gui{0.5\,m$^3$/day}.  Note that a positive value represents
  injection (recharge) and a negative value represents extraction
  (pumping).  This flow rate produces a pore velocity of 2\,m/day for
  the given cross-sectional area and porosity.

\end{itemize}

\begin{tcolorbox}[notebox, title={Note}]
The pore velocity is calculated as $v = Q\,/\,(n \cdot A)
= 0.5\,/\,(0.25 \times 1 \times 1) = 2$\,m/day.  At this velocity,
the advective travel time across the full 100\,m column is 50 days.
The very low dispersivity value (0.01\,m, much smaller than the grid
spacing) is intentional and creates a sharp dissolution front that
closely approximates the ideal advective case.
\end{tcolorbox}

% ============================================================
\section{Running MODFLOW}
% ============================================================

You have now completed the setup of the flow model and are ready to run
MODFLOW\@.  This step provides PHT3D with the steady-state flow field
required to compute advective-dispersive transport.

\begin{itemize}

\item Go to \opath{Parameters $\rightarrow$ Flow} and press the
  \gui{Write} button to write all MODFLOW input files.

\item Press the \gui{Run} button to start the MODFLOW simulation.

\item When the calculation has finished, a window will appear showing the
  last lines of the MODFLOW output file.  Verify that the last line
  confirms a successful simulation.  Because this is a steady-state,
  one-dimensional problem, the simulated head distribution should show a
  uniform gradient along the column from the upstream injection well
  toward the fixed-head downstream boundary.

\item Inspect the flow results by selecting \opath{Results $\rightarrow$
  Flow $\rightarrow$ Head}.  The head field should show a linearly
  decreasing profile from upstream to downstream, consistent with
  uniform flow.

\end{itemize}

% ============================================================
\section{Data Input for Solute Transport}
% ============================================================

The transport parameters control how dissolved species are advected and
dispersed through the model domain.  For this exercise it is assumed
that all chemical components share the same physical transport
properties.

\begin{itemize}

\item Go to \opath{Parameters $\rightarrow$ Transport} and select the
  \gui{Parameter} button.

\item In the dialog that appears, select \gui{ADV} and choose the
  \gui{3rd-order TVD scheme (ULTIMATE)} as the advection solution
  scheme.  Click OK\@.  The TVD scheme is recommended here because the
  very low dispersivity relative to the grid spacing would otherwise
  cause numerical oscillations with a standard finite-difference scheme.

\item Go to \opath{Spatial Attributes $\rightarrow$ MT3DMS
  $\rightarrow$ DSP} and enter a longitudinal dispersivity of
  \gui{0.01\,m}.  Leave all other dispersivity ratios and diffusion
  coefficients at their default values.

\item Go to \opath{Spatial Attributes $\rightarrow$ MT3DMS $\rightarrow$
  BTN $\rightarrow$ btn.11 Effective Porosity} and set the value to
  \gui{0.25}.

\end{itemize}

% ============================================================
\section{PHT3D Reaction Definition}
% ============================================================

In this exercise the geochemical reaction network is defined entirely
through equilibrium reactions between the aqueous solution and the
mineral calcite.  All required species and thermodynamic data are
available in the standard PHREEQC database, so no custom reaction
module is needed.

\begin{itemize}

\item Copy the standard PHREEQC database file into the folder that
  contains all files for this exercise and name it
  \fname{pht3d\_datab.dat}.

\item Go to \opath{Parameters $\rightarrow$ Chemistry} and select the
  \gui{Import database} button.  ORTI3D will read and interpret the
  PHT3D database file.

\item Click on the \gui{Chemical Reaction} button to open the reaction
  definition dialog.

\item In the \gui{Solution} tab, select the aqueous components that will
  be transported.  For this exercise, activate: \gui{C(4)}, \gui{Ca},
  \gui{Cl}, \gui{Na}, \gui{pH}, and \gui{pe}.

\item Click on the \gui{Phases} tab and activate \gui{Calcite} as an
  equilibrium mineral phase.

\end{itemize}

\begin{tcolorbox}[notebox, title={Note}]
The aqueous components C(4) and Ca are the key reactive species.  When
the injected \ce{NaCl} water (which contains no Ca or inorganic carbon)
contacts calcite, the mineral will dissolve to supply both \ce{Ca^{2+}}
and carbonate species to solution.  Na and Cl are conservative tracers
and will travel at the advective pore velocity without reacting.
Including both pH and pe is required for PHREEQC to correctly compute
the carbonate speciation.
\end{tcolorbox}

% ============================================================
\section{Initial and Inflow Concentrations}
% ============================================================

The initial water composition represents a natural groundwater in
equilibrium with calcite.  These concentrations were obtained by running
a PHREEQC equilibrium calculation for a solution in contact with calcite
at $\mathrm{pH} = 9.91$.  The inflow water ($0.001\,\mathrm{mol/L}$
\ce{NaCl}) contains sodium and chloride only, with no Ca or inorganic
carbon.  Both compositions are listed in Tables~\ref{tab:aq} and
\ref{tab:min}.

\begin{table}[ht]
\centering
\caption{Aqueous component concentrations (\molw).}
\label{tab:aq}
\begin{tabular}{@{}L{3.5cm}
                >{\raggedright\arraybackslash}p{4.5cm}
                >{\raggedright\arraybackslash}p{4.5cm}@{}}
\toprule
\textbf{Component}
  & \textbf{Initial (Background)}
  & \textbf{Inflow (Solution 1)} \\
\midrule
pH     & 9.91                      & 7.00 \\
pe     & 4.00                      & 4.00 \\
C(4)   & $1.228 \times 10^{-4}$    & 0.0  \\
Ca     & $1.228 \times 10^{-4}$    & 0.0  \\
Na     & 0.0                       & $1.000 \times 10^{-3}$ \\
Cl     & 0.0                       & $1.000 \times 10^{-3}$ \\
\bottomrule
\end{tabular}
\end{table}

\begin{table}[ht]
\centering
\caption{Mineral concentration (initial background).}
\label{tab:min}
\begin{tabular}{@{}L{5cm}l@{}}
\toprule
\textbf{Mineral}
  & \textbf{Initial concentration (\molvol)} \\
\midrule
Calcite (\ce{CaCO3}) & $1.000 \times 10^{-4}$ \\
\bottomrule
\end{tabular}
\end{table}

\begin{tcolorbox}[notebox, title={Important}]
Mineral concentrations in PHT3D are expressed as moles per litre of
\emph{bulk volume} (\molvol), not per litre of water.  This is
consistent with the convention used in Engesgaard and Kipp (1992) and
in PHT3D Exercise~1 of the course guide.  The initial calcite
concentration of $1.0 \times 10^{-4}$\,\molvol{} is intentionally
small so that it will be entirely consumed by the dissolution front
within the 100-day simulation.
\end{tcolorbox}

\noindent Enter these values into ORTI3D as follows:

\begin{itemize}

\item In the \gui{Solutions} tab of the Chemistry dialog, enter the
  \gui{Background} (initial) water composition using the values from the
  `Initial' column in Table~\ref{tab:aq}.

\item Enter the inflow water composition in the \gui{Solution~1} column
  of the Solutions tab, using the values from the `Inflow' column in
  Table~\ref{tab:aq}.

\item In the \gui{Phases} tab, enter the initial calcite concentration
  from Table~\ref{tab:min} in the \gui{Backgr Mol} field.  Leave the
  \gui{SI} (saturation index) at \gui{0}, indicating that calcite is in
  equilibrium with the initial water composition.

\item To allocate the inflow water composition to the upstream injection
  boundary, go to \opath{Spatial Attributes $\rightarrow$ PHT3D
  $\rightarrow$ PH $\rightarrow$ ph.4 Source Sink} and create a point
  Zone at the same location as the injection well (first cell).  Set the
  Zone Value to \gui{1}.  This instructs PHT3D to use Solution~1 as the
  chemical composition of water entering the column through that cell.

\end{itemize}

% ============================================================
\section{Running PHT3D}
% ============================================================

You are now ready to run the full reactive transport simulation.

\begin{itemize}

\item Go to \opath{Parameters $\rightarrow$ Chemistry} and press the
  \gui{Write} button to generate all PHT3D input files.

\item Press the \gui{Plume} button to start the PHT3D simulation.  A
  command window will appear showing the progress of the calculation.
  The simulation consists of 100 time steps of 1 day each.

\item After the simulation completes, verify that no error messages
  appear and that the run finished at the expected total simulation time
  of 100 days.

\end{itemize}

% ============================================================
\section{Visualization of Simulation Results}
% ============================================================

The primary result to visualize is the spatial distribution of the
calcite dissolution front and the corresponding concentration profiles
of dissolved calcium, inorganic carbon, sodium, chloride, and pH along
the column at selected times.

\subsection{Defining an observation profile}

\begin{itemize}

\item Go to \opath{Spatial Attributes $\rightarrow$ Observation
  $\rightarrow$ OBS $\rightarrow$ obs.1} and use the Zone Tool to draw a
  horizontal line spanning the entire length of the model column.  Name
  this zone \fname{column\_profile}.

\end{itemize}

\subsection{Concentration profiles along the column}

\begin{itemize}

\item In the Results panel, select the time step to visualise by going
  to \opath{Results $\rightarrow$ Aquifer $\rightarrow$ Tstep} and
  selecting, for example, \gui{day 25}.

\item Select a species to display by going to \opath{Results
  $\rightarrow$ Chemistry $\rightarrow$ Species} and choosing \gui{Ca},
  then go to \opath{Results $\rightarrow$ Observation} and select
  \gui{Profile}.  In the dialog, select the profile along
  \fname{column\_profile}.

\item In the following dialog, select \gui{Value} as the averaging
  method.  A concentration profile plot will appear showing the spatial
  distribution of Ca along the column at day 25.

\item Select multiple species simultaneously (Ca, C(4), Na, Cl, pH,
  Calcite) to compare their profiles.  Use the \gui{Export} button to
  save the data to an ASCII file for further post-processing in a
  spreadsheet.

\item Repeat the visualisation for time steps \gui{day 10},
  \gui{day 25}, \gui{day 50}, and \gui{day 75} to observe the
  progressive migration of the dissolution front through the column.

\end{itemize}

\subsection{Space-time evolution as contour maps}

\begin{itemize}

\item Use the contour plot option (\opath{Results $\rightarrow$
  Chemistry $\rightarrow$ Species}) to display the spatial distribution
  of individual species at different times.  Explore how the
  distributions of the conservative tracer \ce{Cl^-} and the reactive
  species Ca and Calcite differ from each other, and consider what this
  implies about the coupling between advective transport and mineral
  dissolution.

\end{itemize}

\noindent If the simulation has been set up correctly, the expected
results should show the following features:

\begin{itemize}

\item Chloride and sodium advance through the column as sharp fronts at
  the advective pore velocity ($\approx 2$\,m/day), acting as
  conservative tracers.

\item Within the zone contacted by the \ce{NaCl} inflow water, calcite
  dissolves and \ce{Ca^{2+}} and C(4) are released into solution.
  Behind the dissolution front, the calcite concentration decreases
  toward zero.

\item Ahead of the dissolution front, the water composition remains
  unchanged, with Ca and C(4) equal to their initial equilibrium values
  ($1.228 \times 10^{-4}$\,\molw) and $\mathrm{pH} = 9.91$.

\item The pH within the zone swept by the \ce{NaCl} water will tend
  toward a lower value reflecting the dissolution equilibrium of the
  residual calcite present, and once calcite is fully consumed, the pH
  will approach that of the injected solution ($\mathrm{pH} = 7$).

\end{itemize}

% ============================================================
\section{Model Variants and Discussion Questions}
% ============================================================

Once you have successfully reproduced the base case simulation, explore
the following model variants and reflect on the results using the
questions below.

\subsection{Effect of initial calcite concentration}

Increase the initial calcite concentration from $1.0 \times 10^{-4}$
to $1.0 \times 10^{-3}$\,\molvol{} and re-run the simulation.

\begin{enumerate}

\item How does the higher calcite concentration affect the position and
  sharpness of the dissolution front after 100 days?  Why?

\item How does the dissolved \ce{Ca^{2+}} concentration profile change?
  What is the maximum dissolved \ce{Ca^{2+}} concentration in the column
  and how does it compare to the equilibrium value?

\item Under what conditions would you expect the dissolution front to
  slow down or stall?

\end{enumerate}

\subsection{Effect of dispersivity}

Reset the initial calcite concentration to $1.0 \times 10^{-4}$\,\molvol.
Then increase the longitudinal dispersivity from 0.01\,m to 1.0\,m and
re-run the simulation.

\begin{enumerate}

\item Compare the concentration profiles of Ca, \ce{Cl^-}, and Calcite
  at day 50 between the low- and high-dispersivity cases.  How does
  dispersivity affect the sharpness of the dissolution front?

\item Does dispersivity affect the \emph{total} amount of calcite
  dissolved after 100 days?  Explain your answer.

\item The course materials show contour plots comparing two different
  simulation scenarios (top and bottom rows in the figure).  Based on
  your experiments, which parameter difference is most likely
  responsible for the contrast between those two rows?

\end{enumerate}

\subsection{Conservative versus reactive fronts}

\begin{enumerate}

\item At day 50, compare the position of the \ce{Cl^-} front
  (conservative tracer) with that of the \ce{Ca^{2+}} front.  Are they
  at the same location?  If not, explain why they might differ.

\item Sketch (or export and annotate) a figure showing the Cl, Ca, C(4),
  and Calcite profiles at days 25, 50, and 75.  Use this to explain the
  concept of a \emph{dissolution front} in your own words.

\item How would the simulation change if the inflow water also contained
  dissolved \ce{CO2} (i.e.,\ C(4) $> 0$ at lower pH)?  Would calcite
  dissolve faster or slower?  Would the dissolution front propagate
  faster or slower than the \ce{Cl^-} front?

\end{enumerate}

\subsection{Connection to \textit{in situ} leaching}

\begin{enumerate}

\item This exercise is conceptually related to \textit{in situ} recovery
  (ISR) of metals, where a lixiviant solution is injected to dissolve a
  target mineral.  Drawing on the results of this exercise, what
  controls the rate at which a dissolution front advances through an ore
  body?

\item If this were a copper-bearing system instead of a calcite system,
  what additional geochemical reactions and complications would you
  expect to encounter?

\end{enumerate}

% ============================================================
\end{document}
% ============================================================